\section{What is a knot statement in first-order arithmetic?}
The objects quantified over in this paper are natural numbers. A diagram is a code $D$, a move is finite data, and a proof that two diagrams represent the same diagrammatic knot is a list $p$. Write
\[
 D\sim_R E\quad\Longleftrightarrow\quad \exists p\;\operatorname{Movie}(D,E,p)=1.
\]
Thus an obstruction has the strong form $I(D)\ne I(E)\Rightarrow
\forall p\,\operatorname{Movie}(D,E,p)=0$. It applies to lists of every length.
A positive construction supplies the code of its witness. An integer, polynomial, group presentation, surface or embedding below always has a specified finite representation.

\paragraph{The arithmetic.}
We use intuitionistic first-order arithmetic with equality, successor axioms, defining equations for primitive recursive (PR) functions, and induction on quantifier-free formulas with parameters. These are precisely the conventions in ProveIt's \texttt{ConciseIPRA/PR.lean} \cite{proveit}. Concretely, pair $a,b$ by $2^a(2b+1)$, with $0$ reserved for the empty list. The projections on positive codes use at most the code many divisions by two. Nest pairs to code lists. Length, lookup with default zero, append, scan and reversal are obtained by loops bounded by the code or length. Arrays of fixed length $l$ with entries $<b$, $b\ge2$, can instead use $l$ base-$b$ digits. Integers are pairs modulo $(u,v)=(u',v')$ iff $u+v'=u'+v$, normalized by truncated subtraction. Rationals have a signed numerator and positive denominator, with order by cross multiplication. Laurent polynomials have an integer low exponent and a list of coefficients, trimmed at the two ends, with a separate zero flag. Addition aligns windows, multiplication is finite convolution. Every algorithm has a default result on a malformed input.

\begin{lemma}[Bounded induction interface]\label{lem:interface}
Bounded conjunction, disjunction, search with a default, double counting of a finite relation and induction over the prefixes of a supplied list can all be used with quantifier-free induction in iPRA.
\end{lemma}
\begin{proof}
For a Boolean $b(i,x)$ iterate $U(0)=1$, $U(k+1)=U(k)b(k,x)$ and $S(0)=0$, $S(k+1)=S(k)+(1-S(k))b(k,x)$. At a failure of $U$ or success of $S$, a simultaneous register stores the first responsible index; a bounded prefix check verifies the answer. Substitution of an arbitrary bounded index into a successful $U$ is proved by induction on the prefix with that index as a parameter. Conversely, to prove $U=1$ from all its instances, the induction formula is $U(k)=1$ under the given premise. Iterate these translations from the innermost bound out. A finite bijection preserves counts by summing its $0,1$ graph by rows and columns. The interchange follows by two inductions on the side lengths and distributivity of addition. On a fixed input list $p$, prove a Boolean assertion about its prefix of length $j$ by induction on $j$ with $p$ free. Checking all smaller instances uses one more bounded scan. No induction on a formula asserting unrestricted existence of a movie is needed.
\end{proof}

\paragraph{Planar input.}
Our primary input is a nonempty finite list of oriented cyclic rational polygonal curves in an oriented plane. There are no collapsed edges, nonincident vertex contacts, overlapped edges or triple points. The only self-contacts and mixed contacts outside consecutive endpoints are transverse double points in edge interiors. A bit at each records which branch lies higher. Crossings are detected by solving two affine linear equations for parameters strictly between zero and one. We allow extra subdivision vertices. A circle without crossings carries its own component record. The same algorithm handles diagrams drawn in a rectangular chart on the sphere; the outside face can be changed by redrawing the chart with a certificate.

A compact version numbers the four darts at each crossing in counterclockwise order, alternately over and under. External edges pair the darts, and an orientation chooses an incoming end on each passage. Store isolated circles separately. Following the opposite dart, then its external edge, traces a component. Using adjacent darts instead gives either smoothing. Components can be counted by marking new orbits and following each for at most the number of darts; for an unoriented pairing graph use its connected components, or count just one direction of each orbit. ProveIt's \texttt{diagram.py} uses the same local convention with the understrand in positions zero and two and checks the face permutation of a connected knot \cite{validator}. A planar drawing supplies a realizability certificate for our inputs, including the relative placements of disconnected pieces. A bare word of crossing visits need not have such a drawing.

\paragraph{Local certificates.}
A redraw supplies an orientation-preserving combinatorial isomorphism of subdivided planar cells containing the drawing, including the unbounded face and the nesting of disjoint curves. One concrete format uses a bounding polygon with oriented triangulations, a correspondence of all oriented triangles preserving incidences, and the subset of edges representing the subdivided drawing. Empty circles are retained. Triangulation validity is checked by pairwise intersection, oriented triangle areas and boundary coverage (the sum of the areas equals the bounding area). A supplied common subdivision permits coordinates to change. A disk replacement cuts out a polygonal disk whose boundary meets the drawing in the prescribed, cyclically ordered ends. With the outside fixed, a local drawing is recorded by its crossings, pairing, boundary order, overpasses and oriented incidences. Accept it by a planar-cell isomorphism to the relevant template up to subdivision. Templates are a simple strand versus a one-crossing curl (I), two disjoint arcs versus an empty bigon with one arc over at both crossings (II), and a triangular exchange of three arcs at distinct height ranks (III). Reflections, rotations, reversals of arc directions and both directions of the replacement are included. In III each pair meets once, with the same pairwise overpass before and after, and the order along each arc reverses. All certificates can use maps to a common triangulated disk with boundary marked. Thus no search for a continuous homeomorphism enters the verifier. On a single crossing the usual sign is fixed by the positive Artin braid generator; mirroring interchanges the signs.

\begin{proposition}[Certificate calculus]
Diagrammatic equivalence is reflexive, symmetric and transitive. Component count is invariant. If a PR observable agrees under a local move and redraw, it agrees along every accepted movie. For any numerical bound $B$, the existence of a movie whose whole code is at most $B$ is decidable in iPRA.
\end{proposition}
\begin{proof}
Use an empty list, reverse the moves and their local isomorphisms, or concatenate lists. Tracing a boundary end through each of the six local tangles gives the same exit before and after; no closed component lies entirely in a moving tangle. Thus the external successor cycles and the isolated components have the same count. For an observable $I$, scan $p$ with prefix invariant $I(D_j)=I(D_0)$, under the fixed assumption of acceptance. For a given $B$ apply the bounded search of Lemma~\ref{lem:interface} to codes $0,\ldots,B$. This proves no bound on the smallest witness as a function of the endpoints.
\end{proof}

\paragraph{Classical interpretation.}
The classical Reidemeister theorem identifies the diagram quotient with tame link types up to ambient isotopy \cite{rei,kauffman}. We use it to identify the conventional topological reading of the diagram theorems. Its full completeness assertion for arbitrarily presented isotopies is external here. The arithmetic conclusions on move lists and the geometric constructions in the next section have independent proofs.
